% General definitions; study-specific recipes are in experimental-study.tex.
% Operational provenance: methodology-notes.md. No empirical results.
\section{Methodology}
\label{sec:methodology}

\subsection{Sparsification interventions}
\label{sec:interventions}

A sparsification intervention specifies a gate-site set $\mathcal{G}$, gate
families and thresholds at those sites, and a pressure method with target set
$\mathcal{P}$. Gate and pressure sites are independently specified: the former
determine where values are set to zero, and the latter where an auxiliary
objective encourages small activations.

We use ReLU and two fixed-threshold gates:
$G_{+,\kappa}(x)=x\mathbf{1}\{x\geq\kappa\}$ and
$G_{\pm,\kappa}(x)=x\mathbf{1}\{|x|\geq\kappa\}$, with $\kappa\geq0$.
One-sided thresholding removes negative values and small positive values;
symmetric thresholding removes small magnitudes while preserving the signs
and values of survivors. At $\kappa=0$, their forward maps are ReLU and the
identity, respectively. Appendix~\ref{app:interventions} specifies boundary
and derivative conventions.

Pressure uses $\mathcal{L}_1$, the equally weighted mean of the mean absolute
activations of each targeted site--layer tensor, taken after its gate when
one is present.
\emph{Naive L1} optimizes $\mathcal{L}_{\mathrm{task}}+\lambda\mathcal{L}_1$.
\emph{Orthogonal L1} (OL1) makes a task-only AdamW update, then adds a pressure
correction preconditioned by the task's second-moment estimates. When this
correction opposes the adaptive task direction, OL1 projects out its
conflicting component. A positive budget caps its norm relative to that task
direction. This controls update geometry without guaranteeing preserved task
loss; Appendix~\ref{app:pressure} defines the procedure.

\subsection{Sparsity and architectural reach}
\label{sec:measurement}

Local sparsity is the fraction of exactly zero entries at a named site.
To weight zeros by the work they affect, we define \emph{model-wide sparsity}
as $\Smodel=N_0/N_{\mathrm{model}}$, reported as $100\Smodel\%$. Here, $N_0$
counts scalar multiplications with at least one zero activation operand in
the declared transformer-block matrix products. $N_{\mathrm{model}}$ counts all valid multiplications
in those operations and the dense language-model head; the head receives no
zero credit. Counts are pooled before division. The percentage refers to this
declared multiplication workload, not all operations or elapsed time. We count
only causally valid pairs in a full, uncached sequence;
Appendix~\ref{app:accounting} gives the exact counters.

The \emph{all-zero reach ceiling},
$\Smax=N_{\mathrm{reach}}(\mathcal{G})/N_{\mathrm{model}}$, counts the fraction
of that workload whose zero operand is guaranteed by setting the selected
gate sites entirely to zero, crediting each affected multiplication once. It depends on the architecture,
site set, and sequence length, but not pressure or a checkpoint. The ceiling
characterizes selected-site reach rather than sparsity achievable at acceptable
quality: observed $\Smodel$ also includes natural zeros outside that reach.
Appendix~\ref{app:ceiling} defines the reach accounting.
